\section{Teoría de juegos}

\subsection{Green Numbers}
Dos jugadores alternan turnos, información completa, sin azar ni empates. \textit{Imparcial} si ambos tienen los mismos movimientos desde todo estado; otro caso es \textit{partidario}. \textit{Normal} el último en mover gana, en \textit{Misère}, pierde.

\textbf{Grafo de estados.} $u\to v$ representa un movimiento. Los terminales son perdedores ($P$); $u$ es ganador ($G$) sii tiene una salida a $P$, y es $P$ sii todas sus salidas van a $G$. Orden topológico inverso. Caracterizar los $P$: propiedad $\mathcal P$: los terminales la cumplen; todo movimiento desde $\mathcal P$ la rompe; desde $\neg\mathcal P$ existe un movimiento a $\mathcal P$.

\textbf{Nim y Grundy}
En Nim, un movimiento reduce una de las pilas $a_1,\dots,a_k$. Sea $X=\bigoplus_i a_i$: la posición es $P$ sii $X=0$. Si $X\ne0$, sea $b$ su bit encendido más significativo; elegir $a_i$ con $b$ encendido y reemplazarla por $a_i\oplus X<a_i$ deja XOR $0$.

Para una posición $s$, $G(s)=\operatorname{mex}\{G(t):s\to t\}$.
Entonces $s$ es $P$ sii $G(s)=0$. \textbf{Sprague--Grundy:} todo juego imparcial finito normal equivale a una pila de Nim de tamaño $G(s)$; para juegos independientes,
$
G_{\rm total}=G(s_1)\oplus\cdots\oplus G(s_k),
$
y el estado global es $P$ sii $G_{\rm total}=0$. 

\textbf{Finales especiales:}  
Si el problema no define los estados perdedores en términos de 'no tener movimientos' (por ejemplo, el objetivo es mover una ficha a una casilla determinada), debemos transformar el juego identificando los nodos que obligan al oponente dejarnos en una posición ganadora. Dichos nodos actuarán como nuestros estados terminales. 

Si quien no puede mover gana, es \textbf{misère} y el XOR de Grundy \textbf{no aplica}.


\textbf{Maximización/minimización.} Si $A$ maximiza $F$ y $B$ la minimiza eligiendo objetos alternadamente: analizar los dos objetos finales para ambos turnos, deducir la condición de intercambio, ordenar con ella y generalizar a $n$ objetos.

\subsection{Green Hackenbush}
Grafo unido al \textit{suelo}; se elimina una arista y desaparece todo lo que pierda conexión con él. Es imparcial normal y las componentes se combinan con XOR.

\textbf{Principio del colon.} Si $G(H_1)=G(H_2)$ con $H_1,H_2$ arboles y  $H_1$ se une al resto del grafo sólo por una articulación, sustituir $H_1$ por $H_2$ no cambia el valor del numero de Grundy de todo el grafo. Esto se reduce a 
$
G(u)=\bigoplus_{v\text{ hijo de }u}(G(v)+1),
$
para arbol enraizado.

\textbf{Principio de fusión.} Fusionar todos los vértices de un ciclo preserva Grundy; sus aristas se vuelven lazos. Unificar el suelo, fusionar ciclos. Reemplazar cada lazo por una hoja y aplicar colon. Lazos en un mismo vértice se cancelan por pares.\\

\algoinfo{ (N+M) \alpha(N) }{ N+M }{Nodos en $[0,N)$}

\inputcode{cpp}{latex_src/game_theory/GreenHackenbush.cpp}{GreenHackenbush.cpp}